\chapter{Conclusions and Remaining Work}
\label{ch:conclusions}

\section{Results established so far}
The work has established a reproducible fixed-mesh phase-field baseline and an operational Abaqus-native pre-refinement chain. The most important current findings are:
\begin{enumerate}
  \item The fixed-mesh Mode-I benchmark reproduces the project reference peak force and displacement closely ($0.757778\,\mathrm{kN}$ at $u=0.005857\,\mathrm{mm}$, error $<0.05\%$), providing a verified quantitative baseline.
  \item Verified direct mesh and contour evidence has been incorporated across the baseline, native-refinement, production validation, and Mode-II studies.
  \item The Pandey--Kumar two-job refinement procedure has been implemented in Python using the Abaqus \texttt{RemeshingRule} and \texttt{adaptiveRemesh} interfaces. Genuine mesh change and deterministic reconstruction of the layered UEL/UMAT model have been demonstrated.
  \item Authoritative serial production solve \texttt{1399632.mmaster02} on the sharp-slit $71{,}320$-element adaptive mesh completed technically to $u=0.010\,\mathrm{mm}$ in $35\,\mathrm{h}\,08\,\mathrm{min}$, resolving a smooth post-peak crack propagation branch ($94.4\%$ load drop) and pure horizontal Mode-I crack path.
  \item Quantitative reproduction of the reference peak force and displacement failed on this mesh ($F_{\text{peak}} = 0.4782\,\mathrm{kN}$ at $u=0.00415\,\mathrm{mm}$, representing $-36.89\%$ and $-29.14\%$ differences vs baseline).
  \item A line-by-line discrepancy audit verified that the UEL formulation, boundary conditions, and seam topology are mathematically correct. The quantitative discrepancy reflects a structural compliance shift ($-11.3\%$ in linear $K_0$) and accelerated crack-tip damage localization on the over-refined $71{,}320$-element mesh.
  \item The nominal $\text{errorTarget}=1.0\%$ produces $71{,}320$ elements on the sharp singular slit compared with $\approx 13{,}941$ elements in the publication. Sizing sensitivity scans confirm that $\text{errorTarget} \approx 2.0\%$ yields $\approx 15{,}396$ elements.
  \item Nonmatching displacement, phase-field and history data can be mapped and ingested by the running UEL, but topology-changing restarts have not yet preserved the relaxed reaction force.
  \item The current shared-state UEL architecture is serial-qualified only. Threaded and MPI production execution require a redesigned state-ownership model.
\end{enumerate}

\section{Open questions}
The immediate scientific questions for Task 5 are:
\begin{enumerate}
  \item What specific adaptivity parameters (error target, pre-analysis state, remeshing region) reconcile the $71{,}320$ vs $13{,}941$ element count and align the adaptive force--displacement curve with the published benchmark?
  \item Can the structural compliance shift and accelerated tip damage localization be fully separated and quantitatively predicted under varying crack-tip mesh densities?
  \item Can a topology-changing restart be formulated so that the fully relaxed receiving state preserves both mechanical force continuity and an internal energy measure?
  \item How should the current shared-state interface be refactored before attempting production multithreading or MPI?
  \item Can the IMFD ABAQUSER visualization tool be integrated into the final workflow and verified against independent ODB field extraction?
\end{enumerate}

\section{Next steps}
Because the quantitative benchmark reproduction gates failed on job \texttt{1399632}, \textbf{Mandatory Thesis Task 5 remains active}. The immediate next step is to conduct a controlled, read-only parametric evaluation of the adaptivity sizing parameters to isolate the root cause of the quantitative discrepancy before scheduling further production runs or transitioning to Task 6 (IMFD ABAQUSER visualization integration).
